\section{Count-localized inversion for the unmodified law}
\label{sec:count-localized-inversion}
The collision count is one of the observed lattice coordinates. Thus a
finite physical-count restriction is exact, not approximate, at every
lattice point whose count lies below the cutoff. This permits the
finite extraction of Section~\ref{sec:finite-count-extraction} to enter
an identity for the original raw density on a central window. The
high-count measure remains in that identity through an explicitly
controlled low-frequency convolution.

Write $\mu_{n,R}=(J_{n,R})_*\nu_R^*$, and let $p_{n,R}$ denote its
mixed density. The density exists by the same countable-word coarea
argument used in Lemma~\ref{lem:finite-graph-density}; the nonnegative
word masses sum to one. For a cutoff $L$, put
\[
 \mu^L_{n,R}=(J_{n,R})_*(\one_{\{N_{n,R}\le L\}}\nu_R^*),
 \qquad \mathsf T^L_{n,R}=\mu_{n,R}-\mu^L_{n,R}.
\]
In this section $E^L,Q^L$ are the exact measures constructed in
Theorem~\ref{thm:finite-count-raw-extraction} for $w=1$, with densities
$e^L,r^L$. We keep the dependence of these choices on $n,L,R$.

\subsection{An exact support property}
\begin{lemma}[Central labels require only a linear count cutoff]
\label{lem:central-count-support}
For $\ell=(k,m)$ with $m\le L$,
\begin{equation}\label{eq:central-count-support}
 p_{n,R}(\ell,t)=p^L_{n,R}(\ell,t)
               =e^L_{n,R}(\ell,t)+r^L_{n,R}(\ell,t)
 \quad\hbox{for almost every }t.
\end{equation}
For $A<\infty$ define the central window
\[
 \mathcal W_{n,R,A}=
 \{(\ell,t)\in\Z^3\times\R:
                  |(\ell,t)-n\bar G_R|\le A\sqrt n\}.
\]
If $\lambda>c_*^{-1}+1$ and $L_n=\lceil\lambda n\rceil$, then
\eqref{eq:central-count-support} holds throughout this window for all
$n\ge\max\{1,A^2\}$, in the mixed almost-everywhere sense.
Moreover, with the constants of
Theorem~\ref{thm:cumulative-return-tail},
\begin{equation}\label{eq:count-tail-for-inversion}
 \|\mathsf T^{L}_{n,R}\|_{\mathrm{TV}}
           \le C e^{an-cL},\qquad
 \sup_\omega|\widehat{\mathsf T^{L}_{n,R}}(\omega)|
           \le C e^{an-cL}.
\end{equation}
\end{lemma}
\begin{proof}
The restriction $N_{n,R}\le L$ is exactly the restriction that the
third lattice coordinate is at most $L$. The measures therefore agree
on each such lattice fiber, and their densities agree almost everywhere
there. The finite extraction supplies the second equality.
On the displayed window the count satisfies
$m\le n/c_*+A\sqrt n\le (c_*^{-1}+1)n$ for $n\ge A^2$.
This proves the asserted coverage. The tail has nonnegative mass equal
to $\nu_R^*(N_{n,R}>L)$; use the cumulative-return tail and then the
basic Fourier bound by total variation. No pointwise density bound for
the high-count measure is deduced from its small mass.
\end{proof}

\subsection{Four terms in an exact inversion}
Choose the cutoff $\chi$ in \eqref{eq:localized-kernel} also to satisfy
$0\le\chi\le1$. Use the proved growing band, namely
\[
 a_n=n^{1/200},\qquad B_n=n^{-99/200},\qquad
 \chi_n(\omega)=\chi(\omega/B_n).
\]
The support has physical radius $2n^{-99/200}$ and rescaled radius
$2n^{1/200}$. The kernel $K_n$ is the mixed inverse transform of
$\chi_n$, with the $(2\pi)^{-4}$ convention of
\eqref{eq:localized-kernel}.

\begin{theorem}[An exact finite extraction on a raw central window]
\label{thm:finite-extraction-central-inversion}
Take $L_n$ as in Lemma~\ref{lem:central-count-support}. On
$\mathcal W_{n,R,A}$, for all sufficiently large $n$,
\begin{align}
 p_{n,R}={}&K_n*\mu_{n,R}
       +(e^{L_n}_{n,R}-K_n*E^{L_n}_{n,R})\notag\\
 &+\mathcal F^{-1}
       [(1-\chi_n)\widehat Q^{L_n}_{n,R}]
       -K_n*\mathsf T^{L_n}_{n,R}.
 \label{eq:count-localized-exact-inversion}
\end{align}
The equality holds almost everywhere. The residual inverse transform
is absolutely convergent, by the proved finite extraction, and
\begin{equation}\label{eq:count-convolution-loss}
 n^2\|K_n*\mathsf T^{L}_{n,R}\|_\infty
          \le C n^{1/50}e^{an-cL}.
\end{equation}
For $L_n=\lceil\lambda n\rceil$ with $\lambda>c_*^{-1}+1$, one also
has, for every $P>0$,
\begin{equation}\label{eq:central-count-Schwartz-tail}
 n^2\sup_{\mathcal W_{n,R,A}}
        |K_n*\mathsf T^{L_n}_{n,R}|\le C_{A,\lambda,P}n^{-P}.
\end{equation}
This last estimate does not require $c\lambda>a$. If that stronger
inequality is imposed, \eqref{eq:count-convolution-loss} additionally
gives exponential smallness on the entire mixed space.
\end{theorem}
\begin{proof}
Theorem~\ref{thm:finite-count-raw-extraction} gives
$\widehat Q^L\in L^1(\T^3\times\R)$, so its inverse is the continuous
representative of the residual density. Splitting that inverse with
$\chi_n$ gives, on the full mixed space,
\[
 p^L=K_n*\mu^L+(e^L-K_n*E^L)
              +\mathcal F^{-1}[(1-\chi_n)\widehat Q^L].
\]
Now substitute $\mu^L=\mu-\mathsf T^L$ and use the exact fiber
agreement of Lemma~\ref{lem:central-count-support} on the window.
Every convolution is defined because its measure is finite and $K_n$
is bounded. This proves \eqref{eq:count-localized-exact-inversion}.
Finally $\|K_n\|_\infty\le CB_n^4$ and
$n^2B_n^4=n^{1/50}$. Combine this with
\eqref{eq:count-tail-for-inversion} to obtain
\eqref{eq:count-convolution-loss}. The power $n^{1/50}$ is retained;
the cutoff error is not declared to be a pointwise tail-density bound.

For the sharper central assertion, the count distance between a point
of the central window and the support of $\mathsf T^{L_n}$ is at least
$(\lambda-c_*^{-1}-1)n$ when $n\ge A^2$. The Schwartz bound in
\eqref{eq:localized-kernel}, and $\|\mathsf T^{L_n}\|_{\mathrm{TV}}\le1$,
therefore give, for every integer $M\ge1$,
\[
 n^2\sup_{\mathcal W_{n,R,A}}|K_n*\mathsf T^{L_n}|
       \le C_{A,\lambda,M} n^{1/50-101M/200}.
\]
Choose $M$ large enough to obtain \eqref{eq:central-count-Schwartz-tail}.
This is decay of the proof kernel across an exactly separated count
support, not decay of the original physical Fourier transform.
\end{proof}

Let
\[
 \mathfrak g_R(y)=\frac{1}{(2\pi)^2\sqrt{\det D_R}}
                    \exp\!\left(-\tfrac12 y^{\mathsf T}D_R^{-1}y\right),
 \qquad y=\frac{(\ell,t)-n\bar G_R}{\sqrt n}.
\]
Uniform positive definiteness was proved in
Theorem~\ref{thm:joint-uniform-nondegeneracy}. Define the actual local
edge correction and finite-band residual integral by
\begin{align*}
 \mathcal E_{n,R,A}^{L}
   &=\mathop{\rm ess\,sup}_{\mathcal W_{n,R,A}}
                   |e^L_{n,R}-K_n*E^L_{n,R}|,\\
 \mathcal C_{n,R}^{L}(B)
   &=\int_{\T^3\times[-B,B]}
               |1-\chi_n(\omega)|\,|\widehat Q^L_{n,R}(\omega)|\dd\omega.
\end{align*}
The first quantity may be infinite until the local density germs are
estimated. Its definition preserves possible cancellation in the exact
edge correction; separate bounds for its two terms also suffice.

\begin{corollary}[A quantitative raw error with the finite branch sum exposed]
\label{cor:finite-extraction-error-budget}
Let $\lambda>c_*^{-1}+1$ and $L_n=\lceil\lambda n\rceil$.
For $P>0$, $B\ge1$, sufficiently large $n$, and every $R\in I$,
the following extended-real inequality holds:
\begin{align}
 &\mathop{\rm ess\,sup}_{\mathcal W_{n,R,A}}
       |n^2p_{n,R}(\ell,t)-\mathfrak g_R(y)|\notag\\
 &\quad\le C n^{-3/280}\sqrt{\log(2+n)}
              +Ce^{-c_0n^{1/100}}
              +C_{A,\lambda,P}n^{-P}\notag\\
 &\qquad\quad+n^2\mathcal E_{n,R,A}^{L_n}
       +(2\pi)^{-4}n^2\mathcal C_{n,R}^{L_n}(B)
       +\frac{n^2}{\pi B}A_2(n,L_n,R,1).
 \label{eq:finite-extraction-error-budget}
\end{align}
Here $C,c_0$ are uniform in the radius; the displayed $A_2$ and edge
terms retain their actual parameter dependence. Thus the raw central
local limit follows if, for some $B=B_n'\ge1$, the last three terms
tend to zero uniformly in $R$ for each fixed $A$.
\end{corollary}
\begin{proof}
In the Fourier integral for $K_n*\mu_{n,R}$, rescale
$v=\sqrt n\,\omega$ and remove the deterministic mean phase.
Theorem~\ref{thm:growing-major-arc} bounds its difference from the
cutoff Gaussian integral by
$C n^{-3/280}\sqrt{\log(2+n)}$ after multiplication by $n^2$.
The complement of the cutoff Gaussian has $|v|\ge n^{1/200}$ and is
bounded by $Ce^{-c_0 n^{1/100}}$ using uniform ellipticity.
These estimates are uniform in the evaluation point, since its Fourier
factor has modulus one.

Apply \eqref{eq:count-localized-exact-inversion}. Its edge term gives
$n^2\mathcal E$, and its high-count convolution gives the third term
in \eqref{eq:finite-extraction-error-budget} by
\eqref{eq:central-count-Schwartz-tail}. Split the absolutely
integrable residual at $|b|=B$. The part inside has precisely the stated
finite-band bound. For $B\ge1$ and sufficiently large $n$, $\chi_n=0$
on $|b|>B$. Equation~\eqref{eq:finite-count-far-tail} bounds this part
by $2(2\pi)^3A_2/B$ before the inverse-transform factor. Since
$(2\pi)^{-4}2(2\pi)^3=1/\pi$, the final constant is as displayed.
This proves the inequality and its sufficient convergence criterion.
\end{proof}

\subsection{Finite bands and weighted observations}
\begin{proposition}[The exact cutoff and a finite-band error budget]
\label{prop:cutoff-finite-band-budget}
For any frequency set $\Omega_n$ of finite volume $V_n$,
\[
 \int_{\Omega_n}|\widehat\mu_{n,R}-\widehat\mu^{L_n}_{n,R}|\dd\omega
       \le C V_n e^{an-cL_n}.
\]
If $V_n\le C e^{\kappa n}$, the choice
\[
 \lambda>\max\{c_*^{-1}+1,(a+\kappa+\delta)/c\},
 \qquad L_n=\lceil\lambda n\rceil,
\]
gives an $O(e^{-\delta n})$ bound and still agrees exactly with the
original raw density on every fixed central window for large $n$.
For a bounded complex weight $w$, the same two statements hold with
an extra factor $\|w\|_\infty$ in the tail estimate. When $w$ is
admissible in Definition~\ref{def:finite-record-weight}, the extraction
and the identity \eqref{eq:count-localized-exact-inversion} also hold
for that weighted measure.
\end{proposition}
\begin{proof}
Integrate \eqref{eq:count-tail-for-inversion} over $\Omega_n$.
The strict inequality in the slope leaves a positive exponential slack.
The support argument is independent of the weight, and the total
variation of the omitted weighted measure is at most its supremum
norm times the omitted probability. Apply
Theorem~\ref{thm:finite-count-raw-extraction} to the weighted finite
measure and repeat the algebra of the inversion identity.
\end{proof}

The last proposition is compatible with the stronger polynomial-weight
truncation estimates in Proposition~\ref{prop:linear-count-budget}.
It is stated here for bounded weights because that is the class used in
the complete finite extraction. For a single initial or actual-return
mark in the inherited BV class, the corresponding central Gaussian
term has the already proved marked estimate, whenever the weight also
satisfies the finite-record subanalytic condition. A finite multiple-time
subanalytic indicator has the structural extraction, but is not thereby
proved to have that marked central estimate. In every case the weight
and its probability denominator remain unchanged.

\paragraph{The remaining long-time estimates.}
The structural finite extraction, its $W^{2,1}$ residual, the exact
central support agreement and the rapidly decaying count correction
are now conclusions. The quantities $\mathcal C_{n,R}^{L_n}(B)$,
$A_2(n,L_n,R,1)$ and $\mathcal E_{n,R,A}^{L_n}$ are the specific remaining
estimates in \eqref{eq:finite-extraction-error-budget}. The first still
contains the intervening annulus, compact nonzero torus frequencies and
the roof frequencies below $B$; its full complement is not bounded by a
one-step function defect. The second has the actual finite-count jet
budget of Proposition~\ref{prop:finite-jet-subtraction}, not a presumed
polynomial growth rate. The third requires local control of singular
values and coefficients on the original central windows. The same
quantitative issues, with the actual indicators, remain in the exact
physical conditioning application. The identity retains the original
raw density throughout; no smoothing or conditioning replaces it.
